\uuid{Pi9a}
\exo7id{1604}
\titre{exo7 1604}
\auteur{barraud}
\organisation{exo7}
\datecreate{2003-09-01}
\isIndication{false}
\isCorrection{false}
\chapitre{Réduction d'endomorphisme, polynôme annulateur}
\sousChapitre{Valeur propre, vecteur propre}
\cpge{mpsi-mp2i.applications-lineaires.b, pcsi.applications-lineaires.b, ptsi.applications-lineaires.b, mp-mpi.reduction-endomorphismes.b, pc.reduction.a, psi.reduction-endomorphismes.a, mp-mpi.reduction-endomorphismes.d, pc.reduction.c, psi.reduction-endomorphismes.c, pt.reduction-endomorphismes.c, mp-mpi.reduction-endomorphismes.g, pc.complements-espaces-vectoriels.d, psi.complements-espaces-vectoriels.d, bcpst2.valeurs-propres-vecteurs-propres.b}
\module{Algèbre}
\niveau{L2}
\difficulte{}

\contenu{
\question{
Soient $u$ et $v$ deux endomorphismes de $E$ qui commutent, c'est \`{a} dire tels que $u\circ
v=v\circ u$. On suppose que $v$ admet $n$ valeurs propres distinctes. Montrer qu'il existe
une base de $E$, form\'{e}e de vecteurs propres communs \`{a} $u$ et \`{a} $v$.

En d\'{e}duire qu'il existe $(a_{0},\ldots,a_{n-1})\in\mathbb{K}^{n}$ tel que
$
 \quad u=a_{0}\mathrm{id}+a_{1}v+\cdots+a_{n-1}v^{n-1}
$
}
}
